\chapter{Modern Algebra}

Modern algebra is concerned with abstract representations of various mathematical objects, and their manipulations and properties.  We will largely be following the material from the book \textit{Algebra} by Michael Artin, as well as some content from Schaum's Outline of Theory and Problems of Abstract Algebra by Frank Ayres, Jr. and Lloyd R. Jaisingh.\\

The reader may be familiar with some of the topics such as groups from courses in physics, but the emphasis here will be to focus on the pure mathematical aspects, with much attention given to theorems and their proofs.  Our treatment will be relatively abbreviated compared to a full course on the subject, but should give the reader enough exposure to feel comfortable with the basics of the subject of abstract algebra.

\section{Groups}

Informally, groups capture the notion of a \textit{symmetry}\index{symmetry}, which can be thought of as a transformation of an object which leaves it apparently unchanged.  For example, consider a square cut from a piece of paper.  If we place the square in front of us flat on the table, then rotate it by one-quarter turn, the square will return to its original appearance.  We can repeat this operation and observe the same result.  We can imagine doing some kind of arithmetic with the square by doing a single rotation, followed by a double rotation, resulting in a sort of addition operation.\\  

To make this more concrete, consider marking each corner of the square by a different letter.  Perhaps the top left corner will be $a$, top right $b$, bottom right $c$ and bottom left $d$.  Let us define a quarter clockwise turn by the symbol $x$, a half clockwise turn $y$, and a three quarter clockwise turn by $z$.  Lastly, let us denote no rotation as $e$, representing the ``identity" symmetry.  We can no specify some particular orientations as follows.  Consider two $x$ operations, which we will call $x^{2}$ (that is, $x \cdot x$).  Now we can see that $x^{2}=y$, as an example.  We now can consider our ``group" to consist of the set of elements $\{x,y,z,e\}$.  Each element represents one of the possible operations.

\section{Rings}

\section{Fields}

\section{Vector Spaces}